\documentclass[fleqn]{llncs}
\usepackage{cite}
%include agda.fmt

%subst keyword a = "\keyword{" a "}"
%subst conid a = "\identifier{" a "}"
%subst varid a = "\identifier{" a "}"
%format (HIDE(x)) = "{}^{\color{gray}" x "}"
%format = = "\kern.15em{=}\kern.15em"
%format -> = →
%format Set = "\constructor{Set}"
%format Set1 = "\constructor{Set_1}"
%format Set₁ = Set1
%format : = "{:}"
%format λ = "\lambda"
%format - = "\text-\kern-.75\underscorespace"
%format . = "\kern.6\underscorespace\text."
%format SEP = "\kern1.55\underscorespace;"
%format ABBRV = "\colorbox[gray]{.9}{\rlap{\ldots}\phantom{\hskip 1.2em}\vphantom{\vrule height 1.1ex}}"
%format tt = "\constructor{tt}"
%format _≡_ = _ "\kern\underscorespace{" ≡ "}" _
%link ≡ _≡_
%format ↦ = "\mathrel{\mapsto}"
%link ↦ _↦_
%format _↦_ = _ "\kern.5\underscorespace{" ↦ "}\kern-\underscorespace" _
%format MAPSTO = ↦
%format /↦ = "\mathrel{\not\mapsto}"
%link /↦ _/↦
%format _/↦ = _ "\kern.5\underscorespace{" /↦ "}"
%format _×_ = _ "{" × "}\kern-\underscorespace" _
%link × _×_
%format _,_ = _ "\kern\underscorespace" , _
%format ⊎ = "\mathrel{\uplus}"
%format _⊎_ = _ "\kern\underscorespace{" ⊎ "}" _
%link ⊎ _⊎_
%format inj₁ = "\constructor{inj_1}"
%format inj₂ = "\constructor{inj_2}"
%format Maybe = "\constructor{Maybe}"
%format nothing = "\constructor{nothing}"
%format just = "\constructor{just}"
%format Bool = "\constructor{Bool}"
%format true = "\constructor{true}"
%format false = "\constructor{false}"
%format if = "\constructor{if}"
%format else = "\constructor{else}"
%format Par = "\constructor{Par}"
%format return = "\constructor{return}"
%format _>>=_ = _ "{" >>= "}\kern-\underscorespace" _
%format fail = "\constructor{fail}"
%format catch = "\constructor{catch}"
%format assert = "\constructor{assert}"
%link assert assert_then_
%format then = "\constructor{then}"
%link then assert_then_
%format assert_then_ = assert _ then "\kern-\underscorespace" _
%link >>= _>>=_
%format Lens = "\constructor{Lens}"
%format Lens.put = Lens . put
%format Lens.get = Lens . get
%format Lens.PutGet = Lens . PutGet
%format Lens.GetPut = Lens . GetPut
%format put' = "\mathit{put}"
%format get' = "\mathit{get}"
%format ↔ = "\scalebox{.75}[1]{$\leftrightarrow$}"
%format _↔_ = _ "{" ↔ "}\kern-\underscorespace" _
%link ↔ _↔_
%format Iso = "\constructor{Iso}"
%format to-from-inverse = to - from - inverse
%format from-to-inverse = from - to - inverse
%format Iso.to = Iso . to
%format Iso.from = Iso . from
%format Iso.to-from-inverse = Iso . to-from-inverse
%format Iso.from-to-inverse = Iso . from-to-inverse
%format iso-lens = iso - lens
%format get-l-a↦b = get - l - a "{\mapsto}" b
%format put-r-b↦b' = put - r - b "{\mapsto}" b'
%format put-l-a-b'↦a' = put - l - a - b' "{\mapsto}" a'
%format get-r-b↦c = get - r - b "{\mapsto}" c
%format List = "\constructor{List}"
%format U = "\constructor{U}"
%format list = "\constructor{list}"
%format μ = "\mu"
%format Pattern = "\constructor{Pattern}"
%format var = "\constructor{var}"
%format k = "\constructor{k}"
%format one = "\constructor{one}"
%format child = "\constructor{child}"
%format left = "\constructor{left}"
%format right = "\constructor{right}"
%format prod = "\constructor{prod}"
%format elem = "\constructor{elem}"
%format Expr = "\constructor{Expr}"
%format BiGUL = "\constructor{BiGUL}"
%format failB = "\constructor{fail}"
%format skip = "\constructor{skip}"
%format replace = "\constructor{replace}"
%format update = "\constructor{update}"
%format rearr = "\constructor{rearr}"
%format dep = "\constructor{dep}"
%format caseS = "\constructor{caseS}"
%format caseV = "\constructor{caseV}"
%format align = "\constructor{align}"
%format BiGULs = "\constructor{BiGULs}"
%format CaseSBranch = "\constructor{CaseSBranch}"
%format CaseSBranchType = "\constructor{CaseSBranchType}"
%format normal = "\constructor{normal}"
%format adaptive = "\constructor{adaptive}"
%format CaseVBranch = "\constructor{CaseVBranch}"
%format replace-lens = "\mathit{replace}" - lens
%format fail-lens = "\mathit{fail}" - lens
%format skip-lens = "\mathit{skip}" - lens
%format caseS-lens = "\mathit{caseS}" - lens
%format put-with-adaptation = put' - "\mathit{with}" - adaptation
%format caseV-lens = "\mathit{caseV}" - lens
%format guarded-return = guarded - "\mathit{return}"
%format DecEq = "\constructor{DecEq}"

\newlength{\underscorespace}
\setlength{\underscorespace}{.8pt}
\newlength{\binderspace}
\setlength{\binderspace}{.8pt}
\newlength{\binderskip}
\setlength{\binderskip}{.8pt}

\newcommand{\keyword}[1]{\mathbf{#1}}
\newcommand{\identifier}[1]{\mathit{#1}}
\newcommand{\constructor}[1]{\mathsf{#1}}

\usepackage[hidelinks]{hyperref}
\renewcommand{\chapterautorefname}{Chapter}
\renewcommand{\sectionautorefname}{Section}
\renewcommand{\subsectionautorefname}{Section}
\renewcommand{\figureautorefname}{Figure}

\usepackage[color=yellow,textsize=footnotesize]{todonotes}

\usepackage{amsmath}

\newcommand{\reason}[1]{\quad\{\,\text{#1}\,\}}

\newcommand{\BiGUL}{\textsc{BiGUL}}
\newcommand{\Agda}{\textsc{Agda}}
\newcommand{\PutGet}{\textsc{PutGet}}
\newcommand{\GetPut}{\textsc{GetPut}}

\usepackage{enumitem}

\begin{document}
\title{A Minimalist Putback-Based Bidirectional Programming Language, Formally Verified}
\author{Hsiang-Shang Ko\kern.05em\inst{2} \and Tao Zan\kern.05em\inst{1} \and Zhenjiang Hu\kern.05em\inst{1,\,2}}
\institute{%
SOKENDAI (The Graduate University for Advanced Studies), Japan\\
{\tt \{taozan,hu\}@@nii.ac.jp}
\and
National Institute of Informatics, Japan\\
{\tt hsiang-shang@@nii.ac.jp}}
\maketitle

\section{Introduction}
\label{sec:introduction}
% Sketch:
%   1. Brief introduce BX, and then Start from Put-based BX
%   2. Talk the problem of Putlenses, BiFluX
%   3. Give our solution
%   4. Show the advantages

% Revise: emphasis on Formally verified,
%         also state Minimal but useful enough since it is a core for BiFluX.

% 1.
%   data replica
%   Internet of Things
%   information explosion
%   Maybe Remove this sentence.
Due to the increase in technology, smart devices for different purpose are developed, thus come into an explosion in the number of devices and information.
The duplication of information among different devices makes it hard to be kept consistent which is an important issue.
% BX for synchronization
Bidirectional transformation (BX)~\cite{GRACE:09} consists a pair of 
$get$ which extracts information from a source to construct an abstract view
and $put$ which propagates updates on view back to the source
, that provides a novel mechanism for maintaining the consistency between two related information (source and view).
% From language approach: bx languages.
Many bidirectional programming languages~\cite{Lenses, Bohannon:06, HuMT08, MHNHT07, Hidaka:10, PachecoCunha:10, Barbosa:2010, Hofmann:2011, FoPP08, HoPW12} are designed
to aid user to write bidirectional transformations,
thus user only needs to write one program that is interpreted in two directions.
The correctness of this BX program is guaranteed by well-behavedness~\cite{Lenses}.
The existing BXs are $get$-based, which let user to write the BX program in the direction of $get$ to describe how to construct the abstract view form source, and system derives a suitable backward $put$ for free.
While in fact there are usually more than one correct $put$ for the $get$, and $get$-based BX only provides a default or limited strategies which hinder BX from real practical usage.
Recently, a new approach called putback-based bidirectional transformation~\cite{FiHP15} was proposed which designs the BX from $put$ instead of $get$
which declares that
it lets user have full control over the $put$ to specify the update policies
and an unique and correct forward transformation $get$ is derived for free.
Typically, a combinator library $putlenses$~\cite{PaHF14} provides a set of basic lens combinators to let user write complex bx programs in the putback-based way,
\textsc{BiFlux}~\cite{PaZH14} is a bidirectional function update language for XML which is a carefully designed putback-based bx language. 
The update program user written using \textsc{BiFluX} in fact is the $put$ function, and the system derives $get$ for free.

% 2. Explain the problem
%    putlenses not fully checked ???
%    biflux not clearly checked, or hard to check ?
%    say: not fully, formally verified.
The pioneer work of putback-based BX gives opportunities to make BX really useful in real applications
as it lets user to have full control over the update strategy.
While the well-behavedness is not fully formal verified for both $putlenses$ and \textsc{BiFluX} which is the most important property for {BX}.
As the forward transformation $get$ is derived for free
and user 
do not need to care how $get$ is derived 
and not have to check the derived $get$ is correct or not manually, 
it is necessary to fully guarantee the program user written is well-behaved by the BX language.
%
% concretely explain how it is in putlenses and biflux
% TODO: bad writing
For example, in \textsc{BiFluX},
the semantics of the basic update operator $replace$ is straight-forward 
and the well-beahvedness for $replace$ is not checked, 
and complex ones like $if$ statement are checked dynamically for the well-behavednss 
while all of them are not formally verified.
Since \textsc{BiFluX} is designed for XML data,
it needs to handle complex regular expressions 
and XML structured data which makes the program hard to check.

% Needs
It turned out that a small but powerful putback-based core language 
which is fully verified to guarantee the well-behavedness is needed 
to serve as the base for all high-level putback-based languages.
Thus we design and implement \BiGUL, a minimalist datatype-generic putback-based language.
% show more of BiGUL
We design \BiGUL\ to be minimal to make it easy for usage and property checking. 
For example, there is no $if$ statement in \BiGUL\ as it can be expressed as a special case of $case$ statement.
While the well-behavedness of \BiGUL\ is fully verified using \Agda\ language.

%   4. Show the advantages
Our contribution can be summarised as follows:
\begin{itemize}[leftmargin=0.3in]
  \item we propose a minimalist putback-based bidirectional language \BiGUL\ as a core language which is fully certified to guarantee the correctness, and served as a base for other putback-based bidirectional languages.
  \item A correct forward transformation $get$ is automatically constructed through \BiGUL\ programs. User concentrates on specifying $put$, our system guarantees the derived $get$ together with the $put$ is well-behaved.
  \item \BiGUL\ has been implemented using Haskell and the well-behavedness has been proved using \Agda\ language and all the source code are available~\footnote{http://www.prg.nii.ac.jp/projects/bigul}.
\end{itemize}

\todo[inline]{Theory construction in terms of \Agda\ programming}

\section{Discussion}
\label{sec:discussion}

\todo[inline]{recursion (the lack thereof)}

\todo[inline]{extensions: for example, iteration, view dependency}

\BiGUL\ is powerful to express a lot of useful BX programs,
and there are extensions can be made to make it support many other useful features.
%
Currently alignment only handles the case that
elements with the same matching key in both source and view are unique % TODO: matching key ?
which indicates the update operations after alignment focus on single source/view element.
While there are situations that more than one elements have the same matching key value,
the result would be a list of elements for the same key value.
\BiGUL\ can be extended with iteration operation on list to hanle this case.
%
% source dependency rearrangement already done.
% view dependency need new feature to hanle it.
Another case is view dependency that means two view elements are denpendent.
For example the catalog of a book dependends on the strcuture and contents of the book.
If the title of contents changes, the catalog shall be changed.
A new feature called view dependency can be added into \BiGUL\ to support it.


\todo[inline]{Haskell implementation}

\todo[inline]{Towards a dependently typed version}

\todo[inline]{stems from work on \textsc{BiFluX}; comparison with putlenses}

\todo[inline]{Theory construction in terms of \Agda\ programming\\
Regarding formalisation: ``[W]e are interested in extending what can be calculated precisely because we are not interested in the calculations themselves. Developing techniques that allow more of a derivation to be calculated allows attention to be focussed on the creative parts, which cannot be calculated.''~\cite{Gibbons-BMF}\\
\cite{UFP-HoTT}}

\bibliographystyle{plain}
\bibliography{ref}

\end{document}
